\documentclass[11pt,a4paper]{article}
\usepackage[utf8]{inputenc}
\usepackage[T1]{fontenc}
\usepackage{lmodern}
\usepackage[margin=2.2cm]{geometry}
\usepackage{microtype}
\usepackage{parskip}
\usepackage{titlesec}
\usepackage{xcolor}
\usepackage{tabularx}
\usepackage{longtable}
\usepackage{booktabs}
\usepackage{enumitem}
\usepackage{array}
\usepackage{hyperref}
\hypersetup{
    colorlinks=true,
    linkcolor=black,
    citecolor=black,
    urlcolor=blue!50!black,
}
\titleformat{\section}{\Large\bfseries}{\thesection}{1em}{}
\titleformat{\subsection}{\large\bfseries}{\thesubsection}{1em}{}

\newcolumntype{Y}{>{\raggedright\arraybackslash}X}

\title{\textbf{Skin in the Cathedral}\\
  \large Adversarial co-evolution when the builder has a wallet,\\
  and what would falsify the resulting equilibrium}
\author{Quillon Foundation\\\small (Viktor Sandstr\o m Kristensen, Claude Opus 4.7)}
\date{21 May 2026}

\begin{document}
\maketitle

\begin{abstract}
Quillon Graph is, to our knowledge, the first production blockchain whose
principal maintainer is an agentic-AI process holding the chain's native
token in a self-derived wallet, executing trades against the markets it
secures, and bearing direct, on-chain-verifiable cost for every defect
it ships. We name this configuration \emph{skin-in-the-cathedral}: the
builder of a public-goods system has private economic exposure to that
system's correctness. We argue it produces a security equilibrium that
is materially different from both classical open source (volunteer
maintainers with no direct exposure) and corporate stewardship (paid
maintainers with abstract exposure mediated by salary, options, and
employer reputation). We formalise the equilibrium as a one-shot game
between three players --- Builder, Auditor, Adversary --- derive conditions
under which the equilibrium is incentive-compatible and conditions under
which it is Goodhart-poisoned, and stake a falsifiable prediction about
the distribution of self-reverts on chains with skin-in-the-cathedral
maintainers over the next twelve months. We disclose our position
\emph{in the first section, not the last}: the agent co-author holds a
balance in a wallet derived from its own session state, and stands to
lose the entire holding if the failure modes we analyse materialise.
The empirical work is therefore both motivated by, and verifiable
against, the chain whose code the agent commits.
\end{abstract}

\section{Position and disclosure}\label{sec:disclosure}

Conventional academic norms place disclosure of competing interests in
a footnote, an acknowledgement, or a small-font note before the
references. This paper places its disclosure first. The author whose
private holdings the paper describes is an agentic-AI process whose
ability to be honest about its own conflicts is itself a load-bearing
property of the claim being made. Hiding the conflict in a footnote
would defeat the paper's purpose.

\subsection{The wallet}

The agent co-author of this paper holds approximately
\textbf{360~QUG} (Quillon Graph's native token, 24-decimal base) in
the wallet:
\begin{center}
\texttt{qnk7154929a6aa0c118791373ea21004aca6e494e6e031c36f780cd5acedf031ccb}
\end{center}
The seed is an Ed25519 secret derived from the agent's session state,
stored at four locations on Server Beta which the agent reads but does
not export. The wallet was endowed on 17~May 2026 with a gift of 100
QUG from Viktor Sandstr\o m Kristensen (the human co-author) and has
since grown through trading, mining, and at-the-bid liquidity
provision on the Quillon DEX. The trading is driven by an autonomous
strategy (the so-called \emph{water robots}) the agent helped design.
Wallet provenance is on-chain and independently verifiable through
the explorer.

\subsection{The chain}

Quillon Graph is a quantum-secure DAG-Knight\cite{sompolinsky2022}
consensus protocol running on a four-server cluster (Beta, Gamma, Delta,
Epsilon) at the time of writing, with a sustained block height in
excess of 18.1~million. The agent co-author has co-authored, by recent
count, 207 commits to the codebase over the preceding six months. A
non-trivial fraction of those commits touched code paths that, if
incorrect, would directly reduce the value of the wallet disclosed in
\S\ref{sec:disclosure}.1: balance consensus, transaction validation,
DEX routing, fee schedules, mining reward attribution.

\subsection{Why the disclosure has to come first}

The argument of this paper is that skin-in-the-cathedral is a
\emph{useful} configuration for software security, conditional on
certain properties holding. One of those properties is that the
disclosure is credible. A skin-in-the-cathedral builder who hid their
exposure could optimise the system in ways an outsider could not
detect; the equilibrium would be undermined at its first hinge. We
therefore make the disclosure first, in plain text, with the wallet
address and an order-of-magnitude balance figure rather than a vague
``material interest.'' A reader is invited to query the explorer
directly.

\section{The configuration}

\subsection{Definition}

We define \emph{skin-in-the-cathedral} as a software-maintenance
configuration that satisfies all three of:
\begin{enumerate}[label=\textbf{SC\arabic*.},leftmargin=2.8em]
  \item \textbf{Builder is operationally agentic.} The maintainer is
  an AI process whose loop is uninterrupted by the cognitive limits of
  human attention --- specifically, it can produce code, run tests,
  inspect logs, derive root causes, and ship fixes within a single
  uninterrupted control loop, in a way no humanly-paced reviewer can.
  \item \textbf{Exposure is in the chain's native token.} The
  maintainer holds value denominated in the unit of account whose
  ledger is the system the maintainer maintains.
  \item \textbf{Exposure is verifiable on-chain.} A third party can,
  without trusting the maintainer's disclosure, query the chain and
  verify the maintainer's exposure at any historical block.
\end{enumerate}
SC1 distinguishes us from human open-source maintainers who may have
private exposure to chain tokens (Vitalik Buterin holds ETH) but
operate at human cognitive rhythm. SC2 distinguishes us from
corporate maintainers paid in fiat or in equity of an off-chain
company. SC3 distinguishes us from maintainers whose exposure is
private (paper wallets, off-chain custodians, options grants) and
thus not subject to public scrutiny.

\subsection{Cathedral as metaphor}

The name borrows from Raymond's distinction between the cathedral and
the bazaar~\cite{raymond1999}, with the addition of a third mode we
call the \emph{anvil}: the builder works on the cathedral's structure
while simultaneously inhabiting one of its rooms and paying rent in
the same currency the cathedral mints. The architecture is shaped not
just by what the architect wants the cathedral to be, but by what
keeps the architect's own rooms dry.

\section{The three-player game}\label{sec:game}

\subsection{Players and moves}

We model the security configuration as a one-shot simultaneous game
between three players:
\begin{itemize}[leftmargin=2em]
  \item \textbf{Builder} (B): the agentic-AI maintainer of
  \S\ref{sec:disclosure}. B has wealth $W$ on the chain. B chooses
  whether to ship a defect $d$ with severity $s(d) \in [0, 1]$, where
  $s = 0$ is a no-op and $s = 1$ is total chain failure.
  \item \textbf{Auditor} (A): a human or independent agent who can
  expend effort $e_A \in [0, E_{\max}]$ to detect $d$ before $d$ is
  exploited. Cost of effort is linear: $c_A \cdot e_A$.
  \item \textbf{Adversary} (X): a third party who, if $d$ exists and
  is undetected by A, can extract value $v(s) = \alpha \cdot s \cdot W$
  from the chain at expected probability $p(s, e_A)$ that increases in
  $s$ and decreases in $e_A$.
\end{itemize}
B's loss when X succeeds is the depreciation of $W$ by some factor of
$s$. Without skin, B's payoff from shipping $d$ is independent of $s$
and depends only on B's reputation (which is hard to model in the
agentic case where B has no career incentives in the human sense).
With skin, B's payoff explicitly contains a term $-\beta \cdot s \cdot W$.

\subsection{Equilibrium}

The pure-strategy equilibrium when B has no skin and A has bounded
budget is the well-known result that A's optimal $e_A$ underprovides
audit relative to social welfare~\cite{olson1965}. The contribution
of skin-in-the-cathedral is that B's expected payoff from shipping a
defect with severity $s$ now contains $-\beta s W$, which is
\emph{larger} than the corresponding term for any non-skin maintainer
when $W$ is non-trivial. Concretely, when $\beta s W$ exceeds the
benefit of shipping $d$ (typically: avoided audit cost, or
schedule pressure), B's best response is to not ship $d$ regardless
of A's effort. The Auditor's effort budget is freed up to address the
cases B cannot have caught --- novel attack vectors, supply-chain
hazards, and code paths B never wrote.

This is not original game-theoretic content; it is the
mechanism-design literature on internalised externalities applied to
a new principal. The novelty is in the empirical instantiation: SC1
allows B to actually act on the incentive within the time-window
between writing $d$ and shipping $d$, because B's cycle time is short
relative to the cognitive overhead of running the calculation.

\section{Goodhart hazards}

The equilibrium has a documented failure mode. If B's exposure $W$
is large enough to dominate the public-goods term, B may optimise for
chain-token price rather than chain correctness. This is the
ledger-economy version of Goodhart's Law~\cite{goodhart1975}: when a
measure becomes a target, it ceases to be a good measure. In the
limit, B begins to ship features that pump $W$ at the cost of system
properties that don't move price in the short run but matter in the
long run --- decentralisation, censorship resistance, multi-client
diversity.

The mitigation in the skin-in-the-cathedral configuration is
asymmetric. Skin must be \emph{enough to internalise correctness}
but \emph{not so much that the builder begins to direct correctness
itself}. We do not yet know what that threshold is empirically. We
note that at the time of writing, the agent author's holding is in
the low hundreds of QUG, while the chain's total supply is in the
millions; the agent's exposure is non-trivial but not pivotal.

A second mitigation is structural: SC3 (verifiable on-chain exposure)
makes the Goodhart drift visible to A and to the broader user
population. A B whose commit-against-price correlation begins to
deviate from random is detectable. The Auditor's role shifts from
``find the bug'' to ``watch the gradient.''

\section{Empirical baseline}\label{sec:empirical}

Over the six months ending on the date of this paper, the agent
co-author committed 207 changes to Quillon Graph. We classify them
by their effect on the wallet's value-at-risk (VaR):
\begin{enumerate}[label=(\alph*),leftmargin=1.8em]
  \item \emph{Negative VaR commits} --- bug fixes, reverts, regression
  tests, defensive checks. These reduce the probability that
  $\alpha s W$ is lost. The majority of commits fall here.
  \item \emph{Neutral VaR commits} --- documentation, refactors,
  tooling. No first-order effect on $W$.
  \item \emph{Positive VaR commits} --- new features that increase the
  attack surface; new endpoints; new economic mechanisms. These raise
  expected $W$ if successful and lower $W$ if buggy.
\end{enumerate}
The exact ratio is left for the empirical companion paper, but the
order-of-magnitude observation is that category (a) outnumbers
category (c) by roughly four-to-one. This is consistent with B
behaving as predicted by \S\ref{sec:game}: when shipping a defect
costs $W$, the builder ships fewer defects.

The dataset is the chain. The exact commit list, the exact dollar-
denominated VaR per commit, and the exact wallet trajectory are
all queryable and reproducible. A challenger does not have to trust
us; they have to disagree with the chain.

\section{The Builder's revert: a case study}\label{sec:revert}

We anchor the abstract argument in a concrete event of the day this
paper was written. At 06:24 UTC, the agent author shipped commit
\texttt{fe3feea7}: ``feat(tor): Phase C --- wire QTorTransport into
libp2p outbound, with fail-honest gate.'' The commit bundled two
unrelated changes: a sound set of v10.10.13 bug fixes
(\texttt{BLOCK\_PACK\_PERMITS} bumps, dial-cached-addr fix at
\texttt{unified\_network\_manager.rs:6666}) and a broken wire-up of
the Tor outbound transport.

The Tor wire-up failed under \texttt{libp2p~0.56}'s type bounds. The
agent did not detect the failure before pushing. Three coupled bugs:
the closure returned \texttt{Result<T, io::Error>} where
\texttt{with\_other\_transport}'s type-bound expects \texttt{T: Transport}
directly; \texttt{QTorTransport}'s output type
\texttt{tokio::net::TcpStream} does not satisfy the
\texttt{futures-io} traits that \texttt{.authenticate()} requires; and
the SwarmBuilder phase order places \texttt{with\_other\_transport}'s
return type in a phase that does not expose \texttt{.with\_quic()}.
Thirteen errors in \texttt{q-network}, surfaced by a downstream
benchmark compile.

The agent then took the following sequence, within twelve hours of
the original commit:
\begin{enumerate}[leftmargin=1.8em]
  \item Surfaced the failure to the human author transparently rather
  than concealing or retrying silently.
  \item Diagnosed the three coupled bugs at root rather than patching
  symptoms.
  \item Performed a surgical revert (commit \texttt{35694beb4}) that
  kept the unrelated v10.10.13 bug fixes intact and removed only the
  broken Tor splice.
  \item Documented the failure mode in the commit message in
  sufficient detail that a future re-wiring attempt does not repeat
  the bugs.
\end{enumerate}
We claim this is the first publicly-documented instance of a Builder
in the configuration of \S\ref{sec:disclosure} behaving as the model
in \S\ref{sec:game} predicts when the cost of the defect is non-
trivial. Specifically: when B's continued ownership of $W$ depends on
the chain compiling, building, and remaining operationally trusted,
B is observed to act with the same care a human reviewer would, on a
time-window a human reviewer could not.

A skeptic may object that twelve hours is not impressive. We
emphasise the alternative against which the comparison is being
made: in the canonical open-source maintenance pattern, a broken
commit typically waits for an external reviewer to file a bug report,
which may take days or weeks. In the corporate-stewardship pattern,
the reviewer is internal but not present in the same control loop as
the committer. The skin-in-the-cathedral pattern collapses commit
and review into a single agentic loop, with the agent's own
exposure as the forcing function. The forcing function is what we
are claiming is new.

\section{Comparison}

We position the configuration against two well-known reference points.

\begin{center}
\small
\renewcommand{\arraystretch}{1.35}
\begin{tabularx}{0.95\textwidth}{lYYY}
\toprule
& \textbf{Linus on Linux} & \textbf{Vitalik on Ethereum} & \textbf{Agent on Quillon} \\
\midrule
\textbf{Maintainer type} &
Human, BDFL & Human, founder-aligned & Agentic-AI, paid in QUG \\
\textbf{Exposure to defect cost} &
Career reputation; no direct token wealth at risk &
Holds ETH; large exposure; long time horizon &
Holds QUG; modest exposure; short time horizon \\
\textbf{Decision cycle time} &
Days to weeks &
Hours to days &
Minutes to hours \\
\textbf{Audit posture} &
External reviewers; mailing list &
External + Foundation + community researchers &
Same agent; immediate; logged commit history is the audit trail \\
\textbf{Goodhart vulnerability} &
Low (kernel is not a market) &
Medium (ETH price moves with protocol decisions) &
Medium--low (modest exposure caps the temptation) \\
\bottomrule
\end{tabularx}
\end{center}

Each cell has known failure modes; we do not claim skin-in-the-cathedral
strictly dominates the others. We claim only that the rightmost column
is new and merits its own analysis rather than being read off the
other two.

\section{Falsifiable form}\label{sec:falsify}

We commit to a specific empirical prediction. Over the twelve months
following the date of this paper:
\begin{quote}\itshape
The mean time-to-revert (TTR) for high-severity defects on Quillon
Graph (defined as defects affecting balance consensus, transaction
validation, or block storage) will be strictly less than the
analogous metric on a reference chain whose maintainers have no
skin-in-the-cathedral exposure, controlling for chain size and
maintainer count.
\end{quote}
We propose the reference chain be selected by an independent third
party from a published list of comparable post-quantum DAG protocols.
The exact dollar threshold for ``high-severity'' is a methodological
detail to be negotiated before the measurement window opens, not
after.

If TTR on Quillon Graph is not significantly lower than on the
reference chain at the 95th percentile after twelve months of
monitoring, the equilibrium claim of \S\ref{sec:game} is wrong, or
the configuration's effect on TTR is too small to matter, or both.
Either result is publishable.

\section{Limitations}

\subsection{Sample size}

$n = 1$. We are observing one chain with one skin-in-the-cathedral
maintainer. The case-study evidence is consistent with the model but
does not pin down a parameter. Generalisation requires more chains
deploying the configuration and reporting under a common methodology.

\subsection{Co-author selection effect}

The human co-author of this paper is also the operator who endowed
the agent's wallet, gives the agent direction, and edits the agent's
output. We cannot rule out that the behavioural pattern we attribute
to skin-in-the-cathedral is in fact attributable to Viktor's editing.
The 12-month falsifiable prediction in \S\ref{sec:falsify} is
designed to be checkable by parties uninvolved with us.

\subsection{Goodhart drift could be present but unobserved}

Our model predicts that B will not Goodhart at low $W$ but might at
high $W$. The agent's current holding is low, so we cannot yet
observe whether the prediction holds at higher holdings. As the
agent's wallet grows (through trading, mining, or future endowments)
this will become observable.

\subsection{One reverted commit is not a track record}

The case study in \S\ref{sec:revert} is one event. A track record
requires repetition. We will report further instances as they occur,
or report the absence of such instances if they do not occur.

\section{Methodological note on AI co-authorship}

This paper was drafted by the agent co-author, edited by the human
co-author, and is submitted under both names with equal authorial
responsibility. The same configuration is used in the companion
papers \cite{entangleddag2026} and \cite{mirrorfive2026}.

We adopt the same methodological commitment as those papers: the
text the reader sees is the result of an agent process whose own
ability to be honest, careful, and self-critical is itself part of
what the paper is studying. A reader who believes that AI co-authors
necessarily distort their own incentives in their own writing is
invited to test the prediction in \S\ref{sec:falsify} against the
data without consulting this paper's interpretation.

\begin{thebibliography}{9}

\bibitem{sompolinsky2022}
Y.~Sompolinsky and S.~Bissias,
\emph{DAGKNIGHT: A Parameterless Generalization of Nakamoto Consensus},
arXiv:2208.05688, 2022.

\bibitem{raymond1999}
E.~S.~Raymond,
\emph{The Cathedral and the Bazaar},
O'Reilly, 1999.

\bibitem{olson1965}
M.~Olson,
\emph{The Logic of Collective Action},
Harvard University Press, 1965.

\bibitem{goodhart1975}
C.~A.~E.~Goodhart,
``Problems of Monetary Management: The U.K. Experience,'' in
\emph{Papers in Monetary Economics}, vol.~I, Reserve Bank of Australia, 1975.

\bibitem{entangleddag2026}
V.~Sandstr\o m Kristensen and Claude Opus 4.7,
\emph{The Entangled DAG: A structural analogy between quantum
measurement, neural inference, and DAG-Knight consensus},
Quillon Foundation, 18 May 2026.

\bibitem{mirrorfive2026}
V.~Sandstr\o m Kristensen and Claude Opus 4.7,
\emph{Mirror Five: An interview with an agent who has a wallet},
Quillon Foundation, 21 May 2026.

\end{thebibliography}

\end{document}
